% Part: second-order-logic
% Chapter: sol-and-set-theory
% Section: power-of-continuum

\documentclass[../../../include/open-logic-section]{subfiles}

\begin{document}

\olfileid{sol}{set}{pow}

\olsection{قوة المتصل}

\begin{explain}
التكميم في الرتبة الثانية يتناول مجموعات !!{domain} الجزئية،
ولا يتناول مجموعات تلك المجموعات الجزئية من !!{domain}.
ولو أردنا هذا التكميم مباشرة لاحتجنا إلى منطق \emph{الرتبة الثالثة}.
فمثلًا، مبرهنة كانتور تنفي وجود دالة !!{injective} من مجموعة
قوى مجموعة إلى المجموعة نفسها. وقد نحاول أن نقول:
«لكل مجموعة~$X$ ولكل مجموعة~$P$، إذا كانت~$P$ مجموعة
قوى~$X$، فليس~$\cardle{P}{X}$».
لكن كون~$P$ مجموعة قوى~$X$ معناه أن !!{element}s~$P$
هي مجموعات~$X$ الجزئية كلها ولا غيرها، فنحتاج إلى عبارة
من نحو~$\lforall[Y][(P(Y) \liff Y \subseteq X)]$.
وموضع الإشكال هو~$P(Y)$؛ إذ ليست هذه !!a{formula} من الرتبة
الثانية: متغير العلاقة الأحادي الموضع، كـ$P$، لا يأخذ في موضع
وسيطه إلا حدًا.

ومع ذلك يمكن \emph{محاكاة} التكميم على مجموعات المجموعات
إذا اتسع المجال لذلك. نستخدم علاقة ثنائية الموضع~$R$،
تربط !!{element}s من !!{domain} بـ!!{element}s من !!{domain}.
فمتى أعطينا~$R$ جمعنا !!{element}s التي يرتبط بها عنصر~$x$
بالعلاقة~$R$. وتكون~$\Setabs{y \in \Domain{M}}{R(x,y)}$
المجموعة التي نرمز إليها بـ$x$.
وعكسًا، لتكن~$Z \subseteq \Pow{\Domain{M}}$ مجموعة من
مجموعات~$\Domain{M}$ الجزئية. فإن كان في~$\Domain{M}$ من
!!{element}s ما لا يقل عددًا عن مجموعات~$Z$، وجدت
علاقة~$R \subseteq \Domain{M}^2$ يكون بها لكل مجموعة رمز:
بوساطة~$R$ يرمز بعض~$x$ إلى كل~$Y \in Z$.
\end{explain}

\begin{defn}
متى كان~$R \subseteq \Domain{M}^2$ قلنا إن
\emph{$x$ يرمز بالعلاقة~$R$} إلى
$\Setabs{y \in \Domain{M}}{R(x, y)}$.
\end{defn}

إذا كان !!a{element}~$x \in \Domain{M}$ يرمز
بالعلاقة~$R$ إلى مجموعة~$Z \subseteq \Domain{M}$،
كانت المجموعة~$Y \subseteq \Domain{M}$ ترمز إلى جملة من المجموعات:
ما ترمز إليه !!{element}s~$Y$.
وبهذا قد تكون~$Y$ رمزًا، بالعلاقة~$R$، لـ$\Pow{X}$.
وشرط ذلك أن يكون لكل~$Z \subseteq X$ عنصر~$x \in Y$
يرمز بالعلاقة~$R$ إلى~$Z$، وأن يكون كل~$x \in Y$
رمزًا بالعلاقة~$R$ لمجموعة~$Z \subseteq X$.

\begin{prop}
!!{formula}
  \[
  \fn{Codes}(x, R, Z) \ident \lforall[y][(Z(y) \liff
    R(x, y))]
  \]
تقرر أن~$s(x)$ يرمز بالعلاقة~$s(R)$ إلى~$s(Z)$.
وأما !!{formula}
\begin{multline*}
  \fn{Pow}(Y, R, X) \ident \\
  \lforall[Z][(Z \subseteq X \lif \lexists[x][(Y(x) \land
    \fn{Codes}(x, R, Z))])] \land {} \\ \lforall[x][(Y(x) \lif
  \lforall[Z][(\fn{Codes}(x, R, Z) \lif Z \subseteq X)])]
\end{multline*}
فتقرر أن~$s(Y)$ ترمز بالعلاقة~$s(R)$ إلى مجموعة قوى~$s(X)$؛
أي إن عناصر~$s(Y)$ ترمز بالعلاقة~$s(R)$ إلى مجموعات~$s(X)$
الجزئية كلها دون غيرها.
\end{prop}

\begin{explain}
بهذا الترميز نعبّر عن قضايا مجموعة القوى، فنجعل التكميم على
رموز المجموعات الجزئية عوضًا عن المجموعات نفسها.
ومن ذلك أن نصوغ مبرهنة كانتور بنفي وجود دالة !!{injective}
من أي مجموعة~$Y$ ترمز بالعلاقة~$R$ إلى مجموعة قوى~$X$، إلى~$X$ نفسها.
\end{explain}

\begin{prop}
الجملة التالية صحيحة منطقيًا:
\begin{multline*}
  \lforall[X][\lforall[Y][\lforall[R][(\fn{Pow}(Y, R, X) \lif \\
      \lnot \lexists[u][(\lforall[x][\lforall[y][((Y(x) \land Y(y) \land
            \eq[u(x)][u(y)]) \lif
            \eq[x][y])]] \land {}\\
        \lforall[x][(Y(x) \lif X(u(x)))])])]]]
\end{multline*}
\end{prop}

\begin{explain}
مجموعة قوى مجموعة !!a{denumerable} هي !!{nonenumerable}؛
فأصليتها تزيد على أصلية أي مجموعة !!{denumerable}، أي
$\aleph_0$. وحجم~$\Pow{\Nat}$ هو «قوة المتصل»، إذ يساوي
حجم مجموعة نقاط المستقيم الحقيقي~$\Real$.
فإن اتسع !!{domain} لترميز مجموعة قوى مجموعة !!a{denumerable}،
  أمكن أن نصف مجموعة بأنها من حجم المتصل بقولنا إنها تكافئ
  عدديًا مجموعة~$Y$ ترمز مجموعة قوى مجموعة~$X$
  من الحجم~$\aleph_0$، بشرط أن يكون لكل مجموعة جزئية من~$X$
  رمز واحد لا غير في~$Y$؛ وهو شرط التفرّد المثبت في
  \(\fn{Cont}\) أدناه. فأما الترميز غير المتفرّد فلا يفيد إلا
  حدًا أدنى لحجم~$Y$. وإن ضاق المجال عن ذلك، فلم توجد
فيه مجموعة جزئية مكافئة عدديًا لـ$\Real$، امتنع وجود علاقة
  ترمز~$\Pow{X}$ أيضًا.
% تصحيح مصدر (0340-I1، 0340-I2): قُيد الحقن بالمجموعة Y وصُرّح بشرط تفرد رموز مجموعة القوى.
\end{explain}

\begin{prop}
إذا كان~$\cardle{\Real}{\Domain{M}}$ كانت !!{formula}
\begin{multline*}
\fn{Cont}(Y) \ident 
\lexists[X][\lexists[R][((\fn{Aleph}_0(X) \land
      \fn{Pow}(Y, R, X)) \land \\ 
      \lforall[x][\lforall[y][((Y(x) \land {} 
      Y(y) \land \lforall[z][R(x,z) \liff R(y,z)]) \lif \eq[x][y])]])]]
\end{multline*}
تعبيرًا عن~$\cardeq{s(Y)}{\Real}$.
\end{prop}

\begin{proof}
تقرر~$\fn{Pow}(Y, R, X)$ أن~$s(Y)$ ترمز
بالعلاقة~$s(R)$ إلى مجموعة قوى~$s(X)$؛
وتقرر~$\fn{Aleph}_0(X)$ أن المجموعة الأصلية قابلة للتعداد.
فـ$s(Y)$ لا يقل حجمها عن قوة المتصل، وقد يزيد عليها إذا
اجتمعت عناصر عدة من~$s(Y)$ ترمز إلى الجزء نفسه من~$X$.
وإنما يمنع المقترن الأخير هذا الاحتمال؛ إذ يشترط أن يكون الربط
بالعلاقة~$s(R)$ بين عناصر~$s(Y)$ ومجموعات~$s(X)$ الجزئية
!!{injective}.
\end{proof}

\begin{prop}
يتحقق~$\cardeq{\Domain{M}}{\Real}$ إذا وفقط إذا
\begin{multline*}
  \Sat{M}{\lexists[X][\lexists[Y][\lexists[R][(\fn{Aleph}_0(X) \land \fn{Pow}(Y, R, X) \land {}\\
      \lforall[x][\lforall[y][((Y(x) \land Y(y) \land
        \lforall[z][R(x,z) \liff R(y,z)]) \lif \eq[x][y])]] \land {}\\
      \lexists[u][(\lforall[x][\lforall[y][(\eq[u(x)][u(y)] \lif \eq[x][y])]] \land {}
        \\\lforall[x][Y(u(x))] \land {}
        \\\lforall[y][(Y(y) \lif \lexists[x][\eq[y][u(x)]])])])]]]}. 
\end{multline*}
\end{prop}

\begin{explain}
مؤدى فرضية المتصل أن حجم المتصل هو أول أصلية
!!{nonenumerable}؛ أي إن حجم~$\Pow{\Nat}$ هو~$\aleph_1$.
\end{explain}

\begin{prop}
صدق فرضية المتصل يكافئ الصحة المنطقية للجملة
\[\fn{CH} \ident
\lforall[X][(\fn{Aleph}_1(X) \liff \fn{Cont}(X))]\]
\end{prop}

ولا يلزم من هذا أن تكون~$\lnot \fn{CH}$ صحيحة منطقيًا إذا
وفقط إذا كذبت فرضية المتصل. فالمجال !!a{enumerable} لا يشتمل
على مجموعة جزئية حجمها~$\aleph_1$، ولا على مجموعة جزئية
من حجم المتصل. لذلك تصدق~$\fn{CH}$ دائمًا في مجال
!!a{enumerable}. غير أن لنفي الفرضية جملة أخرى تكون
صحيحة منطقيًا إذا وفقط إذا كذبت الفرضية:

\begin{prop}
كذب فرضية المتصل يكافئ الصحة المنطقية للجملة
\[\fn{NCH} \ident
\lforall[X][(\fn{Cont}(X) \lif \lexists[Y][(Y \subseteq X \land \lnot
    \fn{Count}(Y) \land \lnot \cardeq{X}{Y})])]\]
\end{prop}

\end{document}
